\section{Best Practices dan Performa CSS}

CSS yang performan dan mudah dirawat memerlukan perhatian pada beberapa aspek: efisiensi selector, minimalisasi reflow/repaint, manajemen ukuran file, dan strategi loading. Dalam proyek skala besar, CSS yang tidak terkelola dapat memperlambat rendering halaman secara signifikan \cite{webdevCss}.

\begin{table}[htbp]
\centering
\begin{tabular}{|P{3.5cm}|P{4cm}|P{4.5cm}|}
\hline
\textbf{Praktik} & \textbf{Cara Implementasi} & \textbf{Dampak} \\
\hline
Animasi GPU-friendly & Gunakan \code{transform}/\code{opacity} & Hindari reflow; animasi mulus \\
\hline
Hindari \code{@import} & Gunakan \code{<link>} atau bundler & Kurangi latency CSS \\
\hline
Critical CSS & Inline CSS above-the-fold & First paint lebih cepat \\
\hline
Minifikasi & Build tool + purge unused & Ukuran file lebih kecil \\
\hline
Selector efisien & Hindari deep nesting & Parsing lebih cepat \\
\hline
\code{will-change} & Hanya saat perlu & GPU prepare; jangan overuse \\
\hline
\end{tabular}
\caption{Best practices performa CSS}
\end{table}

Gunakan DevTools browser (tab Performance dan Coverage) untuk mengidentifikasi bottleneck. Tab Coverage menunjukkan berapa persen CSS yang benar-benar digunakan di halaman---sering kali lebih dari 50\% CSS tidak terpakai. Tool seperti PurgeCSS menghapus kode CSS yang tidak digunakan secara otomatis. Untuk proyek besar, pertimbangkan CSS-in-JS atau CSS Modules untuk memastikan hanya CSS yang relevan yang dimuat \cite{mdnCss}.

\begin{konsep}
\textbf{will-change: Gunakan dengan Bijak.} Properti \code{will-change: transform;} memberi tahu browser untuk menyiapkan optimisasi GPU. Namun, penggunaan berlebihan (\code{* \{ will-change: transform; \}}) justru menghabiskan memori dan menurunkan performa. Gunakan hanya pada elemen yang akan dianimasikan, dan hapus setelah animasi selesai jika memungkinkan \cite{webdevCss}.
\end{konsep}

